\documentclass[UTF8,11pt,a4paper]{ctexart}
\usepackage[margin=23mm]{geometry}
\usepackage{amsmath,amssymb,booktabs,longtable,array}
\usepackage[colorlinks=true,urlcolor=blue,linkcolor=blue,citecolor=blue]{hyperref}
\usepackage{xurl}
\setlength{\emergencystretch}{3em}
\setlength{\tabcolsep}{4pt}
\title{18号赛题：共享归约的物理年龄搜索与缓存}
\author{R103--R109独立实验记录}
\date{2026年10月11日；R102及旧v14报告另行保留}
\begin{document}
\maketitle
\begin{abstract}
本轮仅重排原GET、MULW、ADD、SET并插入COPY，不使用等值别名、重复GET、重算或算术结合重写。
共享点积列内年龄搜索在原336图中的四个窄窗口点积实例减少1236条COPY，4095节点实例从3138降到2196。
缓存及边界beam1将重复列状态搜索归并为三个资源转移，在该大实例只扩展244个状态、命中312次。
完整336输入仍为272合法、47身份模型证书、9预算失败、8启发式失败；新72结构持出组合保留62合法输出，COPY21752到15538。
153项回归和10000小图差分通过。成果是严格声明模型内的质量改进，不是官方Level通过声明。
\end{abstract}

\section{规范与本轮研究范围}
\href{https://www.chaspark.com/\#/races/competitions/1326690549003280384}{官方18题正文}
经根代理重新读取，仍是原DAG重排与COPY插入；窗口$W=2^b$，$b=2,3,4,5$，距离$W$本身合法。
Level仅列0、1、3，分母及平均权重未定义，没有新增GET、spill、重算或等值替换授权。
当前语义保留每个原节点一次、运算符及原有序参数身份，并保留同寄存器RAW/WAR/WAW依赖。
本轮结构筛选进一步要求纯GET/SET/ADD/MULW、唯一且读写集合分离的架构寄存器，各行独立原ADD链使用相同权重序列，全部原节点覆盖。

Sugita等CC2023的距离ISA代码生成包含spill、SSA destruction、move insertion\cite{sugita}，其完备性依赖spill，不能用于证明本题任意COPY-only图可行。
Chen的min-reg研究\cite{chen}启发了按活跃资源寻找调度顺序，但这里的代价是物理年龄和COPY，不是自由寄存器数。
本轮从已有固定列布局没有新质量收益的证据出发，探索列内原就绪指令顺序及物理窗口状态。

\section{算法及缓存合法性}
接受2--4行完整共享归约结构。各行原乘积由同一列共享GET权重与独立私有GET组成，逐列原ADD累积到独立SET。
列内状态保存每行stage、最新累加器年龄、临时GET或乘积年龄，以及权重年龄和已出生标志。
列边界全部原ADD完成，只有各行累加器存活；开始下一列前，没有提前执行未来列的原节点。
首列省略ADD，尾列可执行SET，两个标志独立入键。

若最老存活原身份年龄为$W$，且下一原指令不是它的最后消费，则先COPY刷新；最后消费可以直接执行。
这是合法的固定顺序发射规则，搜索不包含所有可能的COPY时机。
每次最多一个未消费私有GET；默认最多一个未归约乘积，另有两个的实验。
列内按累计COPY费用Dijkstra，进度仅打破同费用平局；列边界有限beam及每列4000状态限额保持有界。
列截断、列边界和贪心COPY均限制搜索，不宣称全DAG最优或无解。

行的未来调度拓扑相同，故按$(stage,accumulator\_age,temporary\_age)$排序仅合并\emph{资源角色的行置换同构}。
不把不同原GET值或原节点输出声明等值。保留真实行映射及路径，命中缓存后逐项恢复当列真实原GET/MULW/ADD/SET ID，并仍使用该节点原有序args。
列内完整键包含所有上述资源状态；列边界键为完整累加器年龄排序及首/尾列标志。
缓存仅在一次调用内，$R,W$、beam和预算参数固定；不同程序或参数不共用缓存。
依赖筛选保证该资源同构和原ADD顺序；非匹配拓扑拒绝该候选，不能判为原题不可行。

若有$K$种实际出现的边界转移、每种搜索扩展至多$S$状态，有限搜索及堆成本依赖$K,S,R,W$，
不能对任意输入无条件称线性。原ID配方应用及最终物化另为$O(N+C)$。
完整列边界的$R$个正年龄互异，候选年龄组合至多$\binom WR$；首/尾标志与beam决定实际缓存键数。
实际4行$W=8$、52到4095节点、边界beam1均只出现3种转移、扩展244状态；这是本族实测，不是所有输入保证。

\section{配对证据、失败及可归因收益}
所有输入、版本、原JSON哈希、失败状态及严格输出均保留。
每个成功输出检查原JSON身份排列、符号及三个本地整数配置（8位有符号、16位无符号、32位有符号回绕），这些位宽不是官方运算宽度。
下表时间均为功能批观察，其他项目可能占CPU；不据此宣称隔离性能因果。

\begin{longtable}{p{.17\textwidth}p{.30\textwidth}p{.29\textwidth}p{.16\textwidth}}
\toprule
实验 & 全部输入与输出 & 共同COPY质量 & 成本观察 \\
\midrule
R103原dot16 & 旧12；直接inline11；4受限拒绝、1预算 & 共同11：974到680；大图1秒未完成 & inline尾部约1秒 \\
R104缓存16 & 旧/inline/pair均12；4受限拒绝 & 4112到2876；4095图3138到2196 & inline p95约276毫秒 \\
R105同250毫秒16 & beam1/2/6均12；4受限拒绝 & 三种beam均2876 & p95约44/89/228毫秒 \\
R106结构持出72直接 & 旧62；新60；12受限拒绝 & 共同60：21716到15502；21改善38同1退化 & 2旧成功被matcher拒绝 \\
R107完整336组合 & 两侧272；47证书9预算8启发 & 194519到192929；5改善267同 & 完整55.667秒 \\
R108同72组合 & 两侧62；9预算1启发 & 21752到15538；21改善41同 & 完整15.754秒 \\
\bottomrule
\end{longtable}

R103没有缓存重复搜索全部列；R104加入资源转移缓存后恢复4095节点图。
R105降低边界beam未损失本16图质量，不能因此断言beam1对所有输入最优。
R106同时变化2/3/4行、2/7/31/113/331/631列、四窗口、原始拓扑排列及MUL/ADD左右有序参数；
原dot固定4行的seed不改变结构，因此不能仅更换它的种子假称新持出。
R106的一个退化为3行2列$W=8$，旧0COPY、列内受限候选1COPY；两个小$W=4$旧成功被严格matcher拒绝。
组合先建立旧incumbent，再尝试新候选，R108保留这三处旧更优/可行结果，未丢失成功输出。

R107的1590条总减少中，\textbf{1236条}来自四个dot实例且选中新年龄候选；
另外354条来自旧pressure候选的有限墙时完成差异，不能归因于本轮算法。
四个原节点规模52/247/1014/4095的$W=8$实例分别28到19、178到124、768到537、3138到2196。
全272成功输出等权COPY/输出13.8634223\%到13.7395767\%，COPY/原41.9057099\%到41.5616020\%。
其余64图没有输出，完整336实际成绩未定义，不能把失败当0或删除后宣称过Level。
同新72图成功62子集输出比例11.4047209\%到9.3227192\%，原分母19.2816359\%到14.4986450\%。
9.32\%只属于这个明确成功子集，不替代全体或官方评分。

R107保存272份完整输出重新读盘：原JSON身份272/272，符号及异初值三整数272/272；
R108保存62份两种重放62/62。并行重放观察耗时分别4.755/6.256秒与1.162/1.505秒，不是隔离吞吐。
根代理另用完全不导入本项目模块的原JSON脚本检查R105的12输出：全部原身份一次、有序args、reg和窗口引用通过，COPY总2876。
R109完整153项回归通过（45.037秒）；10000固定种子两侧9998合法、2严格模型证书，COPY均562，完整harness6.721秒。

\section{必要界、设备和后续余量}
先前独立原JSON下界及根代理审计仍有效：同原336保留全部失败的必要输出分母等权下界11.7792\%，既有272合法子集10.0864\%。
这些仅限声明的原身份/COPY-only模型与合成分布，不能裁定缺少评分分母、平均权重与官方有效数据的Level。
4行、315列、$W=8$的独立归约角色界至少1431COPY；新2196仍不证明最优，也不保证1431可达。
对应输出分母必要界25.8958\%、实际34.9070\%，真正剩余差距约9.011个百分点；不能把必要界提高计作实际COPY减少。

下一优先级是原17个搜索/启发失败和一般共享依赖图：先诊断真实原图，尝试有限窗口恰好满载时由当前指令最后消费携带值的合法串行顺序；未匹配结构不据此证明无解。
另一研究方向是跨列原就绪顺序：允许一行先进入下一列、其他行随后推进，仍保持各原ADD依赖与共享GET单次。
列内空间遗漏跨列布局，已经由3行2列旧0/新1的合法反例证实。
应先采用两列有限状态、独立身份重放和相同原336回放；更大的lookahead可能增加状态和内存成本。
树家族也有差距，但最直接work-first策略此前无新增收益，避免只重复已有顺序或换更容易分布。

实测平台为Windows、CPython3.12.14标准库、Ryzen9 8945HX（16核32线程）、约31.3GiB内存。
本轮没有GPU、LLVM、FPGA或远程付费算力；未测本轮峰值RSS，不能将状态数冒称内存实测。
此前特例工作集510MiB及下界停止245MiB另有旧资源证据，不能泛化成新搜索峰值。
规划最低4核/8GiB/2GiB证据磁盘，建议8核以上/16--32GiB/10GiB磁盘；尚未在最低设备执行本批。
单开发者下一步估计：两列状态候选3--5工作日，独立证明/差分与拓扑扩充2--4日，隔离性能/内存计量1--2日，
规范与提交接口核对1--2日，共7--13日；这是估算，不是已投入工时。
数据/验证器审计与性能计量可并行，原ID合法性模型及固定基线由统一负责人维护。

\section{冻结和复现}
R109冻结UTC为2026-10-11 07:08:16，快照\path{experiments/baseline_age_r109}。
核心规范化SHA-256：
\url{0c3d19724d85e33ce7bc8cfba24feb94c0bb3a0b3770091be50fd8aaa4b28af4}。
R107/108的执行源与冻结仅命令行\path{__main__.py}策略枚举不同，算法文件一致；R109全部源一致。
\path{experiments/r109-age-validation/source-provenance.json}保留逐文件及报告哈希。
R107/R108对照是实际冻结R102源码，而非当前核心同名策略。
默认\texttt{auto}仍为旧v14，显式\texttt{age\_portfolio}保留R102候选并有界尝试年龄搜索。
单键复现可运行\path{scripts/level-stage.ps1}，参数
\texttt{-Strategy age\_portfolio -Baseline phase\_r102 -Round my-age-audit}，每次使用新Round以保留证据。
严格提交输出只含\texttt{block\_id}和\texttt{optimized\_instructions}，官网已知输入接口通过；不会输出stats。
\path{research/level-stage/reduction-age-search.md}保存完整角色、缓存键和失败边界说明。

\begin{thebibliography}{9}
\bibitem{sugita} Shu Sugita et al. A Sound and Complete Algorithm for Code Generation in Distance-Based ISA. CC2023.
\url{https://doi.org/10.1145/3578360.3580263}.
作者全文：\url{https://www.rsg.ci.i.u-tokyo.ac.jp/lab/en/assets/files/cc2023-sugita.pdf}.
\bibitem{chen} Gang Chen. Optimal and Heuristic Min-Reg Scheduling Algorithms for GPU Programs. 作者预印本v4，2023-11-21.
\url{https://arxiv.org/abs/2303.06855v4}.
\end{thebibliography}
\end{document}
